\section{Conclusion}

A consumer monitoring service is judged on the sentences it sends, so we built
and measured the pipeline accordingly. The component that matters is small: a
\HeadParamsRound-parameter pairing head, supervised for free from masks the
benchmarks already publish, deciding which detections across two revisits are
the same object and which unmatched detections are real change. It lifts
detection from the 23--26 F1 band that every non-learned pairing strategy
occupies to \EOnePixelFOne{} on LEVIR-CD, and it carries that to WHU-CD
unchanged at \ETwoPixelFOne, cutting instance count error to \ETwoCountMAE.

Measuring the reports required a metric with no language model inside it.
Change-Fact-Score reduces a report to atomic claims and checks them against
human references and ground-truth masks; it is what shows that a
vision--language model given both images writes fluently while deciding at
chance whether anything changed. Holding grounding fixed and swapping only the
pairing moves report factuality by \EFiveGain{} points, which is the link
between detection quality and product quality stated as a measurement rather
than an assumption.

The graph earns a narrower claim than the one we set out to make. Its
contribution is unmeasurable at crop level for structural reasons, and at
neighbourhood level it is a counting result: handed every crop of a tile, a
language model's scene count degrades badly
(\ESevenStructCountMAE{} instances of error), while aggregating through the
graph first holds it to \ESevenGraphCountMAE{} --- better than the
deterministic template and an order of magnitude better than either
alternative. What that does not yet test is the temporal half of the design,
because both benchmarks carry two acquisitions and the graph is built to
accumulate across many. A dataset with a real revisit history is the experiment
this work now needs, alongside change types beyond buildings and true satellite
ground sampling distances.
